\begin{definitionbox}{Definition --- Topology}
\begin{definition}[Topology]\label{def:topology-on-set}
Let \(X\) be a set.  A \emph{topology} on \(X\) is a collection \(\tau\) of
subsets of \(X\), whose members are called \emph{open sets}, such that:
\begin{enumerate}
  \item \(X\in\tau\) and \(\emptyset\in\tau\);
  \item the union of any family of members of \(\tau\) belongs to \(\tau\);
  \item the intersection of any two members of \(\tau\) belongs to \(\tau\).
\end{enumerate}
\end{definition}
\end{definitionbox}

\begin{remark*}[Standard quantified statement]
\[
\begin{aligned}
&X\in\tau
\land
\emptyset\in\tau
\\
&{}\land
\forall \mathcal{U}\subseteq\tau{:}\;
\bigcup_{U\in\mathcal{U}}U\in\tau
\\
&{}\land
\forall U,V\in\tau{:}\;
U\cap V\in\tau.
\end{aligned}
\]
\end{remark*}

\begin{remark*}[Predicate reading]
\[
U\in\tau
\]
means that \(U\) is an open subset of the topological space whose topology is
\(\tau\).
\end{remark*}

\begin{remark*}[Interpretation]
A topology selects which subsets of \(X\) count as open.  The axioms require
the whole set and the empty set to be open, allow arbitrary unions of open
sets, and preserve openness under finite intersections.
\end{remark*}

\begin{remark*}[Historical note]
This definition is modeled on Definition~3.1 of Willard's
\emph{General Topology} \cite{WillardGeneralTopology}.
\end{remark*}

\LeanFormalizes{def:topology-on-set}{lra-lean}{LRA.VolumeIV.TopologicalSpaces.TopologicalSpace}{TopologyDefinition}{statement}

\DefinitionalRoot

\begin{definitionbox}{Definition --- Topological Space}
\begin{definition}[Topological Space]\label{def:topological-space}
A \emph{topological space} is a pair \((X,\tau)\), where \(X\) is a set and
\(\tau\) is a topology on \(X\).
\end{definition}
\end{definitionbox}

\begin{remark*}[Standard quantified statement]
\[
(X,\tau)
\quad\text{with}\quad
\tau \text{ a topology on } X.
\]
\end{remark*}

\begin{remark*}[Interpretation]
A topological space is a set together with a chosen topology.  The set
supplies the points, and the topology supplies the open subsets used to speak
about neighborhood, limit, and continuity behavior.
\end{remark*}

\begin{remark*}[Historical note]
This definition is modeled on Definition~3.1 of Willard's
\emph{General Topology} \cite{WillardGeneralTopology}.
\end{remark*}

\LeanFormalizes{def:topological-space}{lra-lean}{LRA.VolumeIV.TopologicalSpaces.TopologicalSpace}{TopologicalSpaceDefinition}{statement}

\begin{dependencies}
\begin{itemize}
  \item \hyperref[def:topology-on-set]{Definition: Topology}
\end{itemize}
\end{dependencies}

\begin{definitionbox}{Definition --- Closed Set}
\begin{definition}[Closed Set]\label{def:topological-closed-set}
Let \((X,\tau)\) be a topological space, and let \(E\subseteq X\).  The set
\(E\) is \emph{closed in \(X\)} if \(X\setminus E\) is open in \(X\).
\end{definition}
\end{definitionbox}

\begin{remark*}[Standard quantified statement]
\[
\forall (X,\tau),\ \forall E\subseteq X,\quad
E\text{ is closed in }X
\Longleftrightarrow
X\setminus E\in\tau.
\]
\end{remark*}

\begin{remark*}[Predicate reading]
\[
E\text{ is closed in }X
\]
means that the complement of \(E\), taken relative to the ambient set \(X\),
is open in the topology on \(X\).
\end{remark*}

\begin{remark*}[Interpretation]
Closed sets are the subsets whose outside is open.  In a topological space,
closedness is therefore not an independent primitive: it is the complement-side
language determined by the chosen open sets.
\end{remark*}

\begin{remark*}[Historical note]
This definition is modeled on Definition~3.3 of Willard's
\emph{General Topology} \cite{WillardGeneralTopology}.
\end{remark*}

\LeanFormalizes{def:topological-closed-set}{lra-lean}{LRA.VolumeIV.TopologicalSpaces.TopologicalSpace}{ClosedSetDefinition}{statement}

\begin{dependencies}
\begin{itemize}
  \item \hyperref[def:topological-space]{Definition: Topological Space}
\end{itemize}
\end{dependencies}

\begin{definitionbox}{Definition --- Closed-Set Family of a Topology}
\begin{definition}[Closed-Set Family of a Topology]\label{def:closed-set-family-of-topology}
Let \((X,\tau)\) be a topological space.  The \emph{closed-set family}
associated to \(\tau\) is
\[
\mathcal{F}_{\tau}
=
\{F\subseteq X:X\setminus F\in\tau\}.
\]
\end{definition}
\end{definitionbox}

\begin{remark*}[Standard quantified statement]
\[
\forall X,\ \forall\tau\subseteq\mathcal{P}(X){:}\quad
\mathcal{F}_{\tau}
=
\{F\subseteq X:X\setminus F\in\tau\}.
\]
\end{remark*}

\begin{remark*}[Predicate reading]
\[
F\in\mathcal{F}_{\tau}
\quad\Longleftrightarrow\quad
F\text{ is closed in }(X,\tau).
\]
\end{remark*}

\begin{remark*}[Interpretation]
The topology \(\tau\) lists the open subsets of \(X\).  The closed-set family
\(\mathcal{F}_{\tau}\) lists the complementary subsets: a subset belongs to
\(\mathcal{F}_{\tau}\) exactly when its complement belongs to \(\tau\).
\end{remark*}

\begin{remark*}[Historical note]
This definition names the closed-set family used in Theorem~3.4 of Willard's
\emph{General Topology} \cite{WillardGeneralTopology}.
\end{remark*}

\LeanFormalizes{def:closed-set-family-of-topology}{lra-lean}{LRA.VolumeIV.TopologicalSpaces.TopologicalSpace}{ClosedSetFamilyOfTopology}{statement}

\begin{dependencies}
\begin{itemize}
  \item \hyperref[def:topological-space]{Definition: Topological Space}
  \item \hyperref[def:topological-closed-set]{Definition: Closed Set}
\end{itemize}
\end{dependencies}

\begin{theorembox}{Theorem}
\begin{theorem}[Closed Sets in a Topological Space]
\label{thm:closed-sets-in-topological-space}
Let \((X,\tau)\) be a topological space.  Its closed-set family
\(\mathcal{F}_{\tau}\) contains \(X\) and \(\emptyset\), is closed under
arbitrary intersections, and is closed under finite unions.

\smallskip
\noindent
\hyperref[prf:closed-sets-in-topological-space]{\textit{Go to proof.}}
\end{theorem}
\end{theorembox}

\begin{remark*}[Standard quantified statement]
\[
\forall X,\ \forall\tau\subseteq\mathcal{P}(X){:}\quad
\mathcal{F}_{\tau}\text{ is the closed-set family of }(X,\tau)
\Longrightarrow
X,\emptyset\in\mathcal{F}_{\tau}
\quad\land\quad
\bigcap_{C\in\mathcal{C}}C\in\mathcal{F}_{\tau}
\text{ whenever }\mathcal{C}\subseteq\mathcal{F}_{\tau}
\quad\land\quad
C\cup D\in\mathcal{F}_{\tau}
\text{ whenever }C,D\in\mathcal{F}_{\tau}.
\]
\end{remark*}

\begin{remark*}[Interpretation]
Closed sets obey the dual closure rules to open sets.  Complements turn
arbitrary unions of open sets into arbitrary intersections of closed sets, and
turn finite intersections of open sets into finite unions of closed sets.
\end{remark*}

\begin{remark*}[Historical note]
This theorem is modeled on Theorem~3.4 of Willard's
\emph{General Topology} \cite{WillardGeneralTopology}.
\end{remark*}

\LeanFormalizes{thm:closed-sets-in-topological-space}{lra-lean}{LRA.VolumeIV.TopologicalSpaces.TopologicalSpace}{closed_sets_in_topological_space}{statement}

\begin{dependencies}
\begin{itemize}
  \item \hyperref[def:topological-space]{Definition: Topological Space}
  \item \hyperref[def:topological-closed-set]{Definition: Closed Set}
  \item \hyperref[def:closed-set-family-of-topology]{Definition: Closed-Set Family of a Topology}
\end{itemize}
\end{dependencies}

\begin{theorembox}{Theorem}
\begin{theorem}[Topology from Closed-Set Axioms]
\label{thm:closed-set-characterization-topologies}
Let \(X\) be a set, and let \(\mathcal{C}\) be a collection of subsets of
\(X\).  Suppose that \(X\in\mathcal{C}\), \(\emptyset\in\mathcal{C}\),
arbitrary intersections of members of \(\mathcal{C}\) belong to
\(\mathcal{C}\), and unions of two members of \(\mathcal{C}\) belong to
\(\mathcal{C}\).  Then \(\mathcal{C}\) determines a topology on \(X\) whose
closed sets are exactly the members of \(\mathcal{C}\).

\smallskip
\noindent
\hyperref[prf:closed-set-characterization-topologies]{\textit{Go to proof.}}
\end{theorem}
\end{theorembox}

\begin{remark*}[Standard quantified statement]
\[
\begin{aligned}
&\forall X,\ \forall\mathcal{C}\subseteq\mathcal{P}(X){:}\;
X\in\mathcal{C}
\land
\emptyset\in\mathcal{C}
\land
\forall\mathcal{A}\subseteq\mathcal{C}{:}\;
\bigcap_{A\in\mathcal{A}}A\in\mathcal{C}
\\
&{}\land
\forall C,D\in\mathcal{C}{:}\;
C\cup D\in\mathcal{C}
\\
&\Longrightarrow
\exists\tau{:}\;
\tau\text{ is a topology on }X
\land
\mathcal{C}
=
\{C\subseteq X:X\setminus C\in\tau\}.
\end{aligned}
\]
\end{remark*}

\begin{remark*}[Predicate reading]
\[
C\in\mathcal{C}
\]
means that \(C\) is one of the subsets being prescribed as closed.
\end{remark*}

\begin{remark*}[Interpretation]
The closed-set axioms are sufficient to recover the corresponding topology:
declare a subset open exactly when its complement belongs to
\(\mathcal{C}\).  The stated closed-set closure laws become the topology
axioms after taking complements.
\end{remark*}

\begin{remark*}[Historical note]
This theorem is modeled on Theorem~3.4 of Willard's
\emph{General Topology} \cite{WillardGeneralTopology}.
\end{remark*}

\LeanFormalizes{thm:closed-set-characterization-topologies}{lra-lean}{LRA.VolumeIV.TopologicalSpaces.TopologicalSpace}{topology_from_closed_set_axioms}{statement}

\begin{dependencies}
\begin{itemize}
  \item \hyperref[def:topological-space]{Definition: Topological Space}
  \item \hyperref[thm:closed-sets-in-topological-space]{Theorem: Closed Sets in a Topological Space}
\end{itemize}
\end{dependencies}

\begin{definitionbox}{Definition --- Topological Closure}
\begin{definition}[Topological Closure]\label{def:topological-closure}
Let \(X\) be a topological space, and let \(E\subseteq X\).  The
\emph{closure} of \(E\) in \(X\) is
\[
\overline{E}^{\,X}
=
\mathrm{Cl}_{X}(E)
=
\bigcap\{K\subseteq X:E\subseteq K\text{ and }K\text{ is closed in }X\}.
\]
\end{definition}
\end{definitionbox}

\begin{remark*}[Standard quantified statement]
\[
\begin{gathered}
X\text{ is a topological space},\quad E\subseteq X,\\
\mathrm{Cl}_{X}(E)
=
\bigcap\{K\subseteq X:E\subseteq K\text{ and }K\text{ is closed in }X\}.
\end{gathered}
\]
\end{remark*}

\begin{remark*}[Interpretation]
The closure of \(E\) is the smallest closed subset of the ambient space \(X\)
that contains \(E\).  It keeps every point of \(E\) and adds exactly the points
forced by closedness in the given topology.
\end{remark*}

\begin{remark*}[Notation]
When the ambient space is clear, \(\mathrm{Cl}_{X}(E)\) may be written
as \(\overline{E}\).  The superscript in \(\overline{E}^{\,X}\) records the
ambient space when closure in different spaces must be distinguished.
\end{remark*}

\begin{remark*}[Historical note]
This definition is modeled on Definition~3.5 of Willard's
\emph{General Topology} \cite{WillardGeneralTopology}.
\end{remark*}

\begin{dependencies}
\begin{itemize}
  \item \hyperref[def:topological-space]{Definition: Topological Space}
  \item \hyperref[thm:closed-sets-in-topological-space]{Theorem: Closed Sets in a Topological Space}
\end{itemize}
\end{dependencies}

\begin{lemma}[Closure Monotonicity]
\label{lem:topological-closure-monotone}
Let \(X\) be a topological space, and let \(A,B\subseteq X\).  If
\(A\subseteq B\), then
\[
\mathrm{Cl}_{X}(A)\subseteq\mathrm{Cl}_{X}(B).
\]

\smallskip
\noindent
\hyperref[prf:topological-closure-monotone]{\textit{Go to proof.}}
\end{lemma}

\begin{remark*}[Standard quantified statement]
\[
\begin{gathered}
X\text{ is a topological space},\quad A,B\subseteq X,\\
A\subseteq B
\Longrightarrow
\mathrm{Cl}_{X}(A)\subseteq\mathrm{Cl}_{X}(B).
\end{gathered}
\]
\end{remark*}

\begin{remark*}[Predicate reading]
\[
\begin{aligned}
&A,B\subseteq X,\\
&A\subseteq B
\quad\Longrightarrow\quad
\overline{A}^{\,X}\subseteq\overline{B}^{\,X}.
\end{aligned}
\]
\end{remark*}

\begin{remark*}[Interpretation]
Closure is monotone with respect to inclusion.  Enlarging the original set can
only enlarge, or leave unchanged, the smallest closed set that contains it.
\end{remark*}

\begin{remark*}[Historical note]
This lemma is modeled on Lemma~3.6 of Willard's
\emph{General Topology} \cite{WillardGeneralTopology}.
\end{remark*}

\begin{dependencies}
\begin{itemize}
  \item \hyperref[def:topological-closure]{Definition: Topological Closure}
\end{itemize}
\end{dependencies}
